\section{Introduction}
Depending on the scientific question, molecular simulation may be organized as a single targeted calculation, an exceptionally long trajectory, a multistage protocol, or a coordinated set of related calculations \cite{shaw2010,adorf2018signac}. For such tasks, the computational procedure cannot always be reduced to a fixed prepare--run--analyze sequence, because continuation, restart, and progression to a subsequent stage may depend on results obtained during execution. Adsorption and water-sorption isotherms provide a representative multi-state example: a complete curve is assembled from simulations performed across a series of thermodynamic conditions rather than from a single calculation \cite{dubbeldam2016,wang2023betgab}. In swelling clays, sorption and desorption may further involve structural transitions, path dependence, and hysteresis \cite{wang2025ca}. Consequently, different states may require markedly different sampling efforts; state completion may need to be assessed from evolving statistical behavior; and, in path-dependent campaigns, downstream states may depend on a validated predecessor configuration. Completing an individual simulation run therefore does not necessarily complete a simulation state; the broader challenge is to manage a persistent, stateful simulation campaign.

Many such studies are organized as sequential campaigns rather than as independent calculations. A molecular system may pass through an ordered series of prescribed conditions, and each system--condition combination constitutes a simulation state with its own configuration, sampling progress, analysis status, and provenance. Several run--analyze--decide cycles may be required before that state is accepted. The accepted state can then initialize the next state in the prescribed path. We refer to this class of calculation as a stateful simulation campaign.

Large language model (LLM) agents are increasingly taking on control roles in scientific simulation rather than serving only as code generators or input-file assistants. In MDCrow, the LLM remains in a repeated reasoning--action loop, interpreting the current objective and selecting among expert-designed tools for molecular dynamics setup, execution, and analysis \cite{mdcrow}. Ding et al. propose a more decoupled supervisory architecture built on OpenClaw, in which the general-purpose agent retains global coordination and supervision while task specification, domain procedures, and high-performance computing (HPC) execution are externalized to dedicated skills and infrastructure \cite{openclawchem}. Recent agent-assisted adsorption studies extend this control role to broader sets of related simulations, where LLMs plan dependencies, allocate tasks, and coordinate multi-case grand canonical Monte Carlo (GCMC) calculations rather than acting only within a single simulation \cite{simmof,multiagentgcmc}. Although these systems differ substantially in architecture, they reflect a broader trend in which the LLM is given responsibility not only for individual actions but also for coordinating increasingly long and complex sequences of scientific computation.

As this control role expands, however, a different question becomes increasingly important: does every stage of a long-running simulation campaign require continued LLM involvement? A molecular simulation campaign is not uniformly reasoning-intensive. Some operations become mechanically prescribed once the scientific protocol has been fixed: a rule-based campaign agent can select the next action from approved rules and recorded state, generate state-specific input files from approved templates and parameters, launch or resume simulation segments from selected restarts, initiate analysis, apply validated continuation and acceptance criteria, and archive accepted states. These operations may be numerous, but they are rule-governed and expected to be reproducible. Placing them in an LLM reasoning loop may add little scientific value while increasing context transmission, token use, latency, and exposure to model nondeterminism. The absence of continuous LLM reasoning therefore does not imply the absence of agentic production control.

Other decisions are more state-contingent. The workflow may encounter a convergence pattern whose scientific meaning is unclear, a provenance conflict between archived and downstream states, an unexpected failure mode, or a mismatch between what a numerical counter records and what the campaign state is intended to represent. In such cases, the difficulty is not simply to execute the next command, but to interpret the meaning of the evolving simulation state. LLM reasoning may therefore be most useful at two boundaries: when a scientific objective is first translated into an executable campaign, and when approved workflow rules can no longer safely determine how the campaign should proceed. The key design problem is thus not how much of a simulation workflow an LLM can control, but where flexible interpretation is necessary and how validated actions return to rule-based agentic production.

The need for such a division of labor is especially clear in stateful hybrid grand canonical Monte Carlo--molecular dynamics (GCMC--MD) isotherm campaigns. In swelling clays, molecular exchange is coupled to interlayer hydration, structural expansion or collapse, and sorption hysteresis, so different relative-humidity (RH) states can exhibit substantially different relaxation behavior \cite{wang2025ca,li2019montmorillonite}. Multiple observables may stabilize on different timescales, while sequential adsorption or desorption can introduce predecessor and restart dependencies between states. Extending the calculation across several mineral compositions further multiplies these dependencies. The resulting campaign therefore contains many routine operations suited to rule-based campaign-agent execution, interrupted by occasional states whose convergence, provenance, or failure status requires flexible interpretation. Stateful GCMC--MD therefore provides a representative setting for examining when reasoning-agent intervention is useful and when rule-based agentic control is preferable.

To address this problem, we developed Agent-MD as an architecture that combines a reasoning agent for campaign construction and event-triggered review, a rule-based campaign agent for persistent production, scientific computation and analysis tools, and persistent campaign state and provenance records. The reasoning agent helps translate the research objective into proposed systems, prescribed conditions, and campaign policies for human approval. The rule-based campaign agent then selects and dispatches actions for system preparation, simulation, analysis, continuation, archiving, and progression between simulation states through scientific tools and recorded state. When approved rules can no longer safely resolve a production state, the campaign agent packages a review event and evidence bundle; the reasoning agent proposes a response, and validation occurs before the campaign agent resumes execution. We implement this architecture for stateful GCMC--MD isotherm campaigns through automated mineral preparation, persistent campaign state, adaptive convergence control, provenance-aware state inheritance, and event-triggered review. The framework is evaluated through a five-system campaign across three RH states, an operational census of production control, analysis and recovery of an authentic timestep-accounting incident, and frozen replay of that incident together with a preserved provenance conflict.
